% ---------------------------------------------------------------------------
% Section 2: Related work (target 1,000 words). Adapted from
% docs/RELATED_WORK.md. Two sources discussed there are not cited here because
% they are absent from docs/LITERATURE.md.
% ---------------------------------------------------------------------------

\section{Related work}
\label{sec:related}

\subsection{The closed-form result, and what it settles}

For a marriage market where status is a single scalar and unions must be
hypergamous, the aggregate questions have closed-form answers.
A recent preprint argues the picture \citep{bilderbeek2026}: universal hypergamy is
feasible exactly when one side's status distribution first-order stochastically
dominates the other's, and when it is not, the minimal unmatched mass is the
hypergamy deficit $D = \sup_t\,(G_W(t) - G_M(t))_+$. The finite-population
condition underneath it is Hall's marriage theorem \citep{hall1935}, of which
this deficit is the maximum deficiency; the preprint's contribution is the
continuum form and its application. This is checked against the simulator below
rather than taken on trust. No behavioral strategy changes either fact,
status gains are transfers of rank rather than creations of it, and under
status-ranked preferences the stable matching is unique and positive assortative.
Existence of a stable matching is \citet{gale1962}; uniqueness under a common
ranking is a property of the lattice they generate, set out by \citet{knuth1996},
because when every agent on a side ranks the other side identically the
proposer-optimal and acceptor-optimal matchings coincide.

The consequences for anyone modeling this are severe. Unmatched mass is a
functional of the two marginal status distributions alone, so interventions that
rearrange who meets whom without changing those marginals cannot move it. This
paper does not re-derive the result; it reproduces it against the simulator
(\S\ref{sec:validation-design}) and then changes the constraint.

\subsection{The empirical record, and why a model is the instrument}

The empirical literature is consistent and mostly about marginals.
\citet{becker1973} supplies the complementarity condition: supermodular surplus
implies positive assortative matching, so measured sorting is evidence about the
surplus function. \citet{schwartz2005} trace US educational homogamy from 1940 to
2003, falling then rising. \citet{eika2019} find assortative mating accounts for
a non-negligible share of cross-sectional household income inequality while
\emph{changes} in assortative mating barely move trends: a level effect without a trend effect, and a warning about which one a cross-sectional model
can speak to. \citet{esteve2016} show the educational gender-gap reversal
shifting marriage patterns toward hypogamy with no preference change required.

On the micro side, \citet{fisman2006} randomise speed-dating encounters and find
trait weights differing by sex and selectivity rising with group size;
\citet{lin2013} show trait salience is not separable, which bounds the additive
utility used here. In online settings, \citet{hitsch2010} estimate mate
preferences, feed them into deferred acceptance, and find predicted stable
matches track observed ones closely, concluding that preferences rather than
frictions explain most sorting. \citet{bruch2018} find a desirability hierarchy
consistent across cities, with reply probability falling sharply in the
desirability gap.

Why model rather than estimate? Because every instrument above observes
\emph{realized} matches inside one institution, and none can answer a
counterfactual about a different one. The question here is institutional: what
happens when the list an agent sees is produced by an engagement-optimizing
estimator rather than given. \citet{kanoria2021} and \citet{halaburda2018} take
the same view from the design side, deriving platform welfare from a model of
restricted choice rather than estimating it. \citet{galichon2022} make the
structural position precise for identification, which is why this paper reports
welfare only inside the model. The precedent for computing when the closed form
does not apply is the clearinghouse work of \citet{roth1999peranson}.

\subsection{Rankers, feedback loops, and the visibility constraint}

\citet{he2017} establish the collaborative-filtering architecture used here and
make no claim about markets. \citet{pizzato2010} introduce the reciprocal
recommender and show that accounting for both sides lifts top-ten success
substantially; \citet{tomita2022} describe deploying a matching-theory-based
recommender on a live platform. Reciprocity-aware ranking is solved engineering,
and a paper claiming novelty for its ranker would be claiming novelty for
something built in 2010. \citet{chen2023} go further, running a
matching-model-based ranker on a live platform and measuring realized matches. That is
the strongest constraint on any novelty claim here: same intersection, different
question.

The loop literature is more directly adversarial. \citet{chaney2018} show that
training on data confounded by the system's own recommendations homogenises users
without raising utility. \citet{mansoury2020} show a loop amplifying popularity
bias and hitting minority-group users hardest. \citet{jiang2019} separate the
echo chamber from the filter bubble. \citet{abdollahpouri2019} frames popularity
bias as a fairness problem, and \citet{sonboli2022} show that in multisided
settings fairness criteria for different stakeholders conflict.

Every one of these predicts a loop that \emph{persists or compounds}. The result
here points the other way, and the distinction this paper offers is between a
ranker's deficiency being coverage, which training could fix, and preference
distortion, which it would not. That distinction is not in the loop literature,
and the honest state of it after this study is that it remains a hypothesis: the
trajectory is consistent with it, and neither of the interventions run here
separates training from reaching (\S\ref{sec:sensitivity},
\texttt{docs/PREREGISTRATION.md} E3). A reader should check it hardest of
anything in the paper, and the paper reports that it does not survive being
checked.

\subsection*{Cold start, exploration, and the dynamic market}

Three literatures bound the question and are not engaged elsewhere in this
section. The first is \emph{cold start} in recommender systems: \citet{schein2002}
define the problem as one where recommendations are required for items no user has
rated, and the standard taxonomy separates item from user cold start. The gap
measured here is a cold-start quantity in exactly that sense — the ranker is being
asked to rank agents it has no interactions for — but the literature's usual
remedies (content features, hybrid models, side information) are not what the
trajectory shows happening, which is the ranker accumulating interactions until
coverage saturates.

The second is \emph{exploration}. \citet{li2010} treat the choice of what to show
as a contextual bandit and make optimism about unseen items the mechanism by which
an estimator learns. That is precisely the intervention this paper could not
implement: an explore-then-exploit ranker would deliberately show low-estimate
agents, which would move reach directly rather than through training. None of the
four arms here explores. A reader who expects the bandit framing should know that
the shipped ranker is pure exploitation, that this is why the cold start exists at
all, and that the paper does not test whether exploring would remove it.

The third is \emph{dynamic matching markets}. \citet{akbarpour2020} model a market
where agents arrive and depart stochastically and study the trade-off between
thickness, information and waiting. This paper's market has neither arrivals nor
departures: generational turnover replaces agents but nobody enters an otherwise
empty market and nobody leaves in response to an outcome. That is why every result
here is about a transient that ends, rather than about a steady state that could
persist — in a market with continuous arrivals some share of the population is
always at epoch zero, and the steady-state gap is the object of interest there. It
has not been measured (\texttt{docs/ROADMAP.md} \S1).

Two further results bound what mediation can be credited with.
\citet{roth1999} show truncating a preference list changes outcomes
strategically, so a top-$k$ list is not a neutral convenience.
\citet{ashlagi2017} show that in an unbalanced random market even a slight
imbalance produces an essentially unique stable matching with the long side
matched near-randomly, so some measured asymmetry is competition rather than
curation. Both cautions apply: the truncation here is imposed rather than
chosen, and the market is balanced, so the competition channel is held at zero by
design rather than measured away.

\subsection{Institutional design and the selection problem}

\citet{kanoria2021} show a platform that restricts agents' actions can raise
welfare. \citet{halaburda2018} show a platform can compete by limiting the
choices it offers, because restricting one agent's choice also restricts her
competitors'. \citet{hoppe2009} give the sharpest constraint on the interior
optimum sought here: with two-sided incomplete information the output gain from
assortative matching can be exactly offset by signalling costs, which is what
failure at high signalling cost looks like.

That literature also names the objection this paper had to design around. If an
institution excludes agents, its measured welfare improves for the arithmetic
reason that excluded agents are no longer counted. The comparability rule of
\S\ref{sec:model} is therefore a response to a known objection, not a safeguard
invented here, and both accountings are reported so the artifact is visible.

Two results bound any $W$-side number. \citet{roth1982} establishes that no
procedure is both always stable and always strategy-proof, with the
proposer-optimal matching simultaneously the acceptor-pessimal one;
\citet{knuth1996} supplies the lattice structure that makes those one statement.
Since hypergamy lives in the acceptor's utility, this bias lands where the
paper's question is, and every acceptor-side level reported is a lower bound
across stable matchings.

\subsection{What is new here}

Thirty-eight sources were examined for whether any already does
recommender-mediated matching on multiple traits. Thirty-five do not. Three
require the claim restated: \citet{chen2023}, \citet{pizzato2010} and
\citet{tomita2022} between them rule out any novelty claim resting on the ranker.
The gap that survives is narrower than ``recommender-mediated matching'' and is a
conjunction of three specific absences: no characterisation of unmatched mass
against the visibility constraint, no institutional comparison on a common
welfare basis, and no horizon. What this paper adds is the
recommender-intermediated setting in which visibility is a truncated top-$k$ list
over an estimated multi-attribute utility and no closed form is known. It claims
nothing for its ranker, and no progress on the multidimensional-status,
search-friction or endogenous-status problems the closed-form literature names as
open.

\subsection*{A note on verification}

Reproducing a closed form is verification: it establishes that the machinery is
correct, not that the model describes anything. \citet{windrum2007} and
\citet{fagiolo2007} set out why empirical validation of an agent-based model is a
family of practices rather than a single test, and why calibration alone does not
resolve it. This paper claims verification and declines the validation question.
